% Working section: structural description of the verified October 4 baseline.
% Final adoption and timing results must be bound before submission.
\section{Full-Color Hardware Computation}
\label{sec:hardware}

\subsection{Boundary between local selection and routed acceptance}

The hardware computation implements one color invocation of the integer
transaction algorithm. Its inputs are an immutable placement snapshot,
unit and member metadata, objective data, configuration, and forbidden
transaction identities. Its outputs are selected transactions, updated local
positions, diagnostics, and per-kind counts. The host retains the epoch and
color schedule, native VPR routing and timing analysis, closure budget,
rollback, batch bisection, and feedback refresh. Consequently, selection by
the color computation is a proposal to the routed controller; it does not
establish an improvement in routed critical-path delay.

The native integration enforces this distinction by invoking the shared
integer implementation on a private copy of the position array. It decodes
the selected records into the proposal protocol consumed by the existing
controller. A rejected native closure therefore restores the incumbent at
the host boundary, independently of the local position updates used to
construct a proposal. The native integration is an implementation equivalence
path, not a measurement of host-to-device transfers or board execution.

\begin{figure}[t]
\centering
\small
\begin{tabular}{@{}c@{}}
Host snapshot and forbidden identities\\
$\downarrow$\\
\fbox{Eligibility and five-kind enumeration}\\
$\downarrow$\\
\fbox{Trial legality and integer evaluation}\\
$\downarrow$\\
\fbox{External filtering and identity encoding}\\
$\downarrow$\\
\fbox{Per-group best, deduplication, conflicts}\\
$\downarrow$\\
\fbox{Batch validation, local commit, statistics}\\
$\downarrow$\\
Host proposal decoding and native routed oracle
\end{tabular}
\caption{Functional order within one color invocation and its host boundary.
The boxes do not denote concurrent streaming stages or an achieved
initiation interval.}
\label{fig:hardware-color}
\end{figure}

\subsection{Explicit identities and conflict sets}

The top-level computation enumerates moves, swaps, macro swaps, and chains
of two or three units. Descriptors retain their indices throughout
evaluation, including unchanged and rejected candidates with zero score.
Canonical identities encode the operation kind, complete nested unit members,
and trial coordinates. External forbidden identities are compared across
their respective pools using complete identity contents. Pool offsets or
host unit numbers alone are insufficient to identify the same transaction.

For each descriptor, evaluation derives touched-net ranges and the sorted,
deduplicated union of old and trial occupied sites. Both sides of a move
belong to its conflict set: checking only destination sites would omit
dependencies created by vacating a location. Explicit offsets carry the
descriptor's ranges through evaluation and encoding. Selection applies the
stable ranking, per-group choice, identity deduplication, and conflict checks
before any local position commit.

The commit operation uses two passes. The first validates the entire selected
batch, including record ranges, known block identifiers, and repeated blocks.
Only after these checks succeed does the second pass write positions.
Evaluation itself leaves positions unchanged. This contract assumes canonical
input backing and pairwise disjoint scratch buffers; scratch from a failed
invocation is discarded. It provides algorithmic batch validation, without
claiming a fault-tolerant memory transaction mechanism.

\subsection{Word interface and capacity}

The synthesized top exposes explicit arrays with 8-, 32-, and 64-bit element
types, 51 logical AXI master bundles, and an AXI-Lite control interface.
Configuration occupies indices 0--35. Fourteen result words report stage and
status, enumeration and movement counts, and the five accepted-kind counts.
The complete C/RTL comparison checks all 30 writable ports, including scratch
arrays and positions, rather than checking only a final objective value.

Logical master bundles do not specify physical HBM channel assignments.
The current out-of-context (OOC) implementation does not establish platform
arbitration, bandwidth, host-device transfer cost, or application throughput.
Likewise, the fixture's \texttt{depth=4096} interface declarations are not proof
that arbitrary benchmark instances fit. The deployment contract must also
account for input counts, checked ranges, scratch capacities, host allocation,
and the eventual memory mapping.

\subsection{Preserving semantics while changing scheduling}

The timing experiments change scheduling and function sharing while retaining
integer widths and transaction semantics. An explicit C++ temporary is not
sufficient evidence of a cycle boundary: HLS may retain a combinational
same-state bypass. Each boundary experiment therefore inspects the generated
RTL, identifies the capture and consuming states, and binds that observation
to the exact generated source. A separate experiment inlines trial-location
decoding to reduce a shared address-selection path; its structural check
examines the resulting module and address logic instead of claiming a new
register boundary.

The physical flow reads a 4\,ns clock constraint before synthesis for
\texttt{xcu280-fsvh2892-2L-e}, then places, optimizes, and routes the generated
design. Acceptance requires the same candidate to pass full functional
regression and final setup, hold, pulse-width, internal-constraint, routing,
and DRC checks. The process must actually exit successfully, with frozen
source and output hashes verified. HLS frequency estimates and intermediate
placement estimates cannot substitute for the final routed report. Because
this OOC flow supplies no input/output delays, even a passing result would
establish internal timing only; physical U280 performance requires separate
platform integration and execution evidence.
